\section*{الفصل \ref{ch:b3:rovibrational-spectroscopy} --- المطيافية الدورانية والاهتزازية}

\begin{solution}{exo:b3:rovibrational-spectroscopy:1}
$I = h/(8\pi^2cB) = 6.62607 \times 10^{-34}/(8\pi^2 \times 2.99792 \times 10^{10} \times
1.9225) = \qty{1.4561e-46}{kg.m^2}$. $\mu = \qty{6.85621}{u} = \qty{1.13850e-26}{kg}$، ومنه
$r_0 = \sqrt{I/\mu} = \qty{113.09}{pm}$.
\end{solution}

\begin{solution}{exo:b3:rovibrational-spectroscopy:2}
$2B_0(J+1)$: 20.88 و41.76 و62.64 و\qty{83.52}{cm^{-1}}. $J_{\max} = \sqrt{kT/2hcB} -
\frac12 = \sqrt{208.5/20.88} - 0.5 = 2.7$: فالمستوى $J = 3$ هو الأكثر
إشغالًا، والخط الأشد هو $4 \leftarrow 3$ عند \qty{83.5}{cm^{-1}}.
\end{solution}

\begin{solution}{exo:b3:rovibrational-spectroscopy:3}
يحتاج الامتصاص في الأمواج الميكروية إلى ثنائي قطب دائم: \ce{HCl} و\ce{H2O} فقط.
ويحتاج رامان الدوراني إلى قابلية استقطاب غير متماثلة الاتجاه: \ce{H2} و\ce{HCl} و\ce{N2}
و\ce{CO2} و\ce{H2O}؛ لا الخذروفان الكرويان \ce{CH4} و\ce{SF6}.
\end{solution}

\begin{solution}{exo:b3:rovibrational-spectroscopy:4}
$N_1/N_0 = \eu^{-hc\tilde\nu/kT}$: عند \qty{300}{K}، $\eu^{-13.84} = \num{9.8e-7}$؛ وعند
\qty{1000}{K}، $\eu^{-4.15} = 0.016$. فلا تصبح \omterm{def:b3:rovibrational-spectroscopy:anharmonic}{الحزمة الساخنة} ملحوظة (واحدًا في المئة من
الأساسية) إلا قرب \qty{1000}{K}.
\end{solution}

\begin{solution}{exo:b3:rovibrational-spectroscopy:5}
$\tilde\omega_e - 2\tilde\omega_ex_e = 2885.3$ و$2\tilde\omega_e - 6\tilde\omega_ex_e =
5665.0$: بطرح ضعف الأولى من الثانية، $-2\tilde\omega_ex_e = -105.6$،
ومنه $\tilde\omega_ex_e = \qty{52.8}{cm^{-1}}$ و$\tilde\omega_e = \qty{2990.9}{cm^{-1}}$.
\end{solution}

\begin{solution}{exo:b3:rovibrational-spectroscopy:6}
$D_e = 2990.95^2/(4 \times 52.82) = \qty{42341}{cm^{-1}}$؛ $G(0) = 1495.47 - 13.20 =
\qty{1482.3}{cm^{-1}}$؛ $D_0 = \qty{40859}{cm^{-1}} = \qty{488.8}{kJ/mol}$، مقابل القيمة الحقيقية
\qty{35760}{cm^{-1}} $= \qty{427.8}{kJ/mol}$: يبالغ نموذج مورس بنسبة 14\,\%.
\end{solution}

\begin{solution}{exo:b3:rovibrational-spectroscopy:7}
$R(0) - P(2) = 6B_0 = 62.64$: $B_0 = \qty{10.440}{cm^{-1}}$. $R(1) - P(1) = 6B_1 = 60.80$:
$B_1 = \qty{10.133}{cm^{-1}}$. $\alpha_e = \qty{0.307}{cm^{-1}}$. وتعطي $R(0) = \tilde\nu_0 + 2B_1$
القيمة $\tilde\nu_0 = \qty{2885.31}{cm^{-1}}$ (و$P(1) = \tilde\nu_0 - 2B_0$ توافقها).
\end{solution}

\begin{solution}{exo:b3:rovibrational-spectroscopy:8}
$\mu_H = \qty{0.97959}{u}$، $\mu_D = \qty{1.90441}{u}$، $\mu_H/\mu_D = 0.51438$. $B_0(\ce{DCl})
= 10.440 \times 0.51438 = \qty{5.370}{cm^{-1}}$. $\tilde\omega_e = 2990.95 \times \sqrt{0.51438}
= 2145.1$، $\tilde\omega_ex_e = 52.82 \times 0.51438 = 27.17$: $\tilde\nu_0 = 2145.1 - 54.3 =
\qty{2090.8}{cm^{-1}}$.
\end{solution}

\begin{solution}{exo:b3:rovibrational-spectroscopy:9}
الانزياحات $4B_0(J + \frac32)$: 11.94 و19.90 و\qty{27.85}{cm^{-1}} (من $J = 0$ و1 و2). 
ونواتا \ce{^{14}N} بوزونان متطابقان سبينهما 1: للمستويات ذات $J$ الزوجي ست
حالات لسبين النوى، وللمستويات ذات $J$ الفردي ثلاث، فتكون الخطوط المنطلقة من $J$ زوجي أشد
بمرتين من جاراتها.
\end{solution}

\begin{solution}{exo:b3:rovibrational-spectroscopy:10}
\omterm{def:b3:quantum-model-systems:rotor}{الدوّار الصلب} يضع الخطوط عند $\nu_1$ و$2\nu_1$ و$3\nu_1$: $230542.40$ 
و\qty{345813.61}{MHz}، أي فوق القيمتين المقيستين بمقدار 4.4 و\qty{17.6}{MHz}. ومع
$\nu_{J+1\leftarrow J} = 2B(J+1) - 4D(J+1)^3$: $\nu_1 = 2B - 4D$، $\nu_3 = 6B - 108D$، ومنه
$3\nu_1 - \nu_3 = 96D$: $D = \qty{0.1835}{MHz}$ و$B = (\nu_1 + 4D)/2 =
\qty{57635.97}{MHz}$. للتحقق: $\nu_2 = 4B - 32D = \qty{230538.0}{MHz}$.
\end{solution}

\begin{solution}{exo:b3:rovibrational-spectroscopy:11}
تنعدم $\Delta G(v + \frac12) = \tilde\omega_e - 2\tilde\omega_ex_e(v+1)$ قرب $v =
\tilde\omega_e/2\tilde\omega_ex_e - 1$؛ ومجموع متتالية حسابية من $n$ حدًا
تنزل من $\tilde\omega_e - 2\tilde\omega_ex_e$ إلى ما يقارب 0 هو نحو $n\tilde\omega_e/2 \approx
\tilde\omega_e^2/4\tilde\omega_ex_e - \tilde\omega_e/2 \approx D_e - G(0)$. ومن أجل \ce{HCl}، مجموع الفجوات الثماني والعشرين
($v = 0$ إلى 27) هو \qty{40857}{cm^{-1}}، مقابل $D_e - G(0) = \qty{40859}{cm^{-1}}$.
\end{solution}

\begin{solution}{exo:b3:rovibrational-spectroscopy:12}
بوجود قيم $J$ الزوجية وحدها، تنطلق الانتقالات $J \to J + 2$ من $J = 0, 2, 4, \dots$: الانزياحات
$6B, 14B, 22B, \dots$، والتباعد $8B$. وللجزيء \ce{CO2} مركز تناظر وليس له
ثنائي قطب دائم: فلا طيف له في الأمواج الميكروية.
\end{solution}

\begin{solution}{pb:b3:rovibrational-spectroscopy:1}
\textbf{1.} $B_0 = \nu/2 = \qty{57.6356}{GHz}$؛ $B_0/c = \qty{1.92252}{cm^{-1}}$.
\textbf{2.} $\mu = 12 \times 15.99491/27.99491 = \qty{6.85621}{u} = \qty{1.13850e-26}{kg}$.
\textbf{3.} $I = h/8\pi^2B_0 = \qty{1.4560e-46}{kg.m^2}$.
\textbf{4.} $r_0 = \sqrt{I/\mu} = \qty{113.09}{pm}$.
\textbf{5.} يزيد $r_0$ على $r_e = \qty{112.83}{pm}$ بمقدار \qty{0.26}{pm}: فالثابت $B_0$ يأخذ متوسط $1/r^2$
على اهتزاز النقطة الصفرية في بئر لاتوافقية، يقضي فيها الجزيء وقتًا أطول
عند المسافات الكبيرة.
\textbf{6.} $2\nu_{1\leftarrow0} = \qty{230.5424}{GHz}$؛ والخط المقيس أدنى بمقدار \qty{4.4}{MHz}:
وهذا هو التشوه الطارد المركزي.
\textbf{7.} $0$ و3.845 و\qty{11.535}{cm^{-1}}، أي $E/k = 0$ و5.53 و\qty{16.60}{K}.
\textbf{8.} $N_2/N_1 = \frac53\exp[-(16.60 - 5.53)\,\mathrm K/T]$.
\textbf{9.} $\ln(\frac53/0.85) = 0.673 = 11.06/T$: $T = \qty{16}{K}$.
\textbf{10.} $J_{\max} = \sqrt{16.4/(2 \times 1.43878 \times 1.9225)} - 0.5 = 1.2$: $J = 1$.
\textbf{11.} يوافق الكمُّ الاهتزازي $hc\tilde\nu/k \approx \qty{3080}{K}$: فعند
\qty{16}{K} لا يكون أي جزيء مثارًا اهتزازيًا، فلا تصدر السحابة أي ضوء
اهتزازي.
\textbf{12.} لا يحتاج إلا إلى \qty{5.5}{K} من الإثارة، فهو مثار حتى في أبرد
غاز، وCO وفير وله عزم ثنائي قطب.
\textbf{13.} $\mu = 13.00335 \times 15.99491/28.99826 = \qty{7.17241}{u}$.
\textbf{14.} $B \propto 1/\mu$: $115.2712 \times 6.85621/7.17241 = \qty{110.1893}{GHz}$.
\textbf{15.} التنبؤ أدنى بمقدار \qty{12}{MHz} (0.011\,\%) من الخط المقيس: $B_0 = B_e - \frac12\alpha_e$، 
والحد الاهتزازي--الدوراني $\alpha_e$ يتغير بقوة مختلفة من $\mu$؛ فبنية
التوازن هي نفسها، أما المتوسط عند النقطة الصفرية فليس كذلك.
\textbf{16.} $98.9/1.1 = 90$.
\textbf{17.} خط \ce{^{12}C^{16}O} مشبع (كثيف ضوئيًا): لم يعد
يشتد مع كمية الغاز.
\textbf{18.} خط \ce{^{13}C^{16}O} النادر يبقى متناسبًا مع كمية الغاز، 
وخط \ce{^{12}C^{16}O} يعطي درجة الحرارة: فهما معًا يقيسان الأمرين كليهما.
\textbf{19.} $\tilde\nu_0 = 2169.81 - 2 \times 13.288 = \qty{2143.24}{cm^{-1}}$، أي \qty{4.666}{\micro m}.
\textbf{20.} $2\tilde\omega_e - 6\tilde\omega_ex_e = \qty{4259.89}{cm^{-1}}$.
\textbf{21.} $G(0) = \frac12\tilde\omega_e - \frac14\tilde\omega_ex_e = \qty{1081.58}{cm^{-1}}$.
\textbf{22.} $R(0) - P(1) = 2(B_0 + B_1) \approx 4B \approx \qty{7.7}{cm^{-1}}$.
\textbf{23.} \ce{^{13}C^{16}O}: \omterm{def:b3:quantum-model-systems:oscillator}{كتلة مختزلة} أثقل، فتردد أدنى.
\textbf{24.} \textbf{$r_0(\ce{CO}) = \qty{113.1}{pm}$}، من تردد راديوي واحد.
\end{solution}
